%% ======================================================================
%%  Index (pp. 108-109)
%% ======================================================================
\clearpage
\folio{108}%
\global\let\currentseccode\empty
\global\let\pagefirstsec\empty
\renewcommand{\chaptertitle}{}%
\def\currentsectitle{}%
\thispagestyle{emptyhead}%
\phantomsection
\pdfbookmark[0]{Index}{fm:index}%
\label{fm:index}%
\addtocontents{toc}{\string\fml{INDEX}{\thepage}{fm:index}}%
\begin{center}{\Large Index}\end{center}%
\vspace{0.5\baselineskip}%
\noindent\small The notation V.2 stands for Chapter V, §2.%
\vspace{0.8\baselineskip}%
\par

\begin{multicols}{2}\small
\setlength{\parindent}{0pt}%
\setlength{\parskip}{0.18\baselineskip}%
\raggedright

abelian extension \hyperref[sec:6.1]{VI.1}

algebraic over a field \hyperref[sec:2.3]{II.3}

algebraic extension \hyperref[sec:2.3]{II.3}

algebraically closed field \hyperref[sec:2.3]{II.3}

associates in a ring \hyperref[sec:1.1]{I.1}

automorphism, Frobenius \hyperref[sec:6.1]{VI.1} and \hyperref[sec:6.3]{VI.3}

\medskip
base, canonical \hyperref[sec:1.4]{I.4}

base of a modulo \hyperref[sec:1.4]{I.4}

bases, dual for the trace \hyperref[sec:2.7]{II.7}

Bezout, identity of \hyperref[sec:1.1]{I.1}

\medskip
characteristic of a field \hyperref[sec:1.7]{I.7}

characteristic polynomial \hyperref[sec:2.6]{II.6}

classes, ideal \hyperref[sec:3.4]{III.4}

closure, integral \hyperref[sec:2.1]{II.1}

conjugate elements or fields \hyperref[sec:2.4]{II.4}

conjugate prime ideals \hyperref[sec:6.2]{VI.2}

cyclic extension \hyperref[sec:6.1]{VI.1}

cyclotomic field, polynomial \hyperref[sec:2.9]{II.9}

cyclotomic extension \hyperref[sec:6.1]{VI.1}

\medskip
decomposition group \hyperref[sec:6.2]{VI.2}

Dedekind ring \hyperref[sec:3.4]{III.4}

degree, residual \hyperref[sec:5.2]{V.2}

dependence, equation of integral \hyperref[sec:2.1]{II.1}

descent, infinite \hyperref[sec:1.2]{I.2}

diophantine equation \hyperref[sec:1.2]{I.2}

discriminant \hyperref[sec:2.7]{II.7}

discriminant, absolute (of a number field) \hyperref[sec:2.8]{II.8}

discriminant ideal \hyperref[sec:2.7]{II.7} and \hyperref[sec:5.3]{V.3}

domain, fundamental \hyperref[sec:4.1]{IV.1}

\medskip
Eisenstein's irreducibility criterion \hyperref[sec:2.9]{II.9}

equation, diophantine \hyperref[sec:1.2]{I.2}

\medskip
{\itshape The index, I copied from old Vladivostok telephone directory}\\
{\itshape Tom Lehrer}

\medskip
equation of integral dependence \hyperref[sec:2.1]{II.1}

equation of Pell-Fermat \hyperref[sec:4.6]{IV.6}

Euler's criterion \hyperref[sec:5.5]{V.5}

Euler's $\varphi$-function \hyperref[sec:1.3]{I.3}

extension, abelian \hyperref[sec:6.1]{VI.1}

extension, algebraic \hyperref[sec:2.3]{II.3}

extension, cyclic \hyperref[sec:6.1]{VI.1}

extension, cyclotomic \hyperref[sec:6.1]{VI.1}

extension, Galois \hyperref[sec:6.1]{VI.1}

extension, quadratic \hyperref[sec:6.1]{VI.1}

\medskip
Fermat's equation \hyperref[sec:1.2]{I.2}

field, cubic \hyperref[sec:2.8]{II.8}

field, cyclotomic \hyperref[sec:2.9]{II.9}

field, number (or algebraic number field) \hyperref[sec:2.3]{II.3}

field, quadratic \hyperref[sec:2.5]{II.5}

finite type, module of \hyperref[sec:1.4]{I.4}

fractions, ring of \hyperref[sec:5.1]{V.1}

fractional ideal \hyperref[sec:3.3]{III.3}

free module \hyperref[sec:1.4]{I.4}

Frobenius automorphism \hyperref[sec:6.1]{VI.1} and \hyperref[sec:6.3]{VI.3}

fundamental units \hyperref[sec:4.4]{IV.4} and \hyperref[sec:4.6]{IV.6}

\medskip
Galois extension \hyperref[sec:6.1]{VI.1}

Galois group \hyperref[sec:6.1]{VI.1}

Gaussian integers \hyperref[sec:5.6]{V.6}

Gaussian sum \hyperref[sec:5.5]{V.5}

greatest common divisor or gcd \hyperref[sec:1.1]{I.1}

group, decomposition \hyperref[sec:6.2]{VI.2}

group, Galois \hyperref[sec:6.1]{VI.1}

group, inertia \hyperref[sec:6.2]{VI.2}

\medskip
ideal, discriminant \hyperref[sec:2.7]{II.7} and \hyperref[sec:5.4]{V.4}

ideal, integral \hyperref[sec:3.3]{III.3}

ideal, fractional \hyperref[sec:3.3]{III.3}

\medskip
identity, Bezout's \hyperref[sec:1.1]{I.1}

imaginary quadratic field \hyperref[sec:2.5]{II.5}

imbedding of a number field, canonical \hyperref[sec:4.2]{IV.2}

index of ramification \hyperref[sec:5.2]{V.2}

inertia group \hyperref[sec:6.2]{VI.2}

integer or integral over a ring \hyperref[sec:2.1]{II.1}

integer of a number field \hyperref[sec:2.8]{II.8}

integer, Gaussian \hyperref[sec:5.6]{V.6}

integral closure \hyperref[sec:2.1]{II.1}

integral dependence equation \hyperref[sec:2.1]{II.1}

integral ideal \hyperref[sec:3.3]{III.3}

integrally closed ring \hyperref[sec:2.2]{II.2}

invariant factors \hyperref[sec:1.5]{I.5}

\medskip
least common multiple or lcm \hyperref[sec:1.1]{I.1}

lattice in $\mathbf{R}^{n}$ \hyperref[sec:4.1]{IV.1}

law, quadratic reciprocity \hyperref[sec:5.5]{V.5}

Legendre symbol \hyperref[sec:5.5]{V.5}

\medskip
minimal polynomial \hyperref[sec:2.3]{II.3}

module of finite type \hyperref[sec:1.4]{I.4}

module, free \hyperref[sec:1.4]{I.4}

module, Noetherian \hyperref[sec:3.1]{III.1}

module, torsion-free \hyperref[sec:1.5]{I.5}

\medskip
Noetherian ring or module \hyperref[sec:3.1]{III.1}

non-residue, quadratic \hyperref[sec:5.4]{V.4}

norm \hyperref[sec:2.6]{II.6}

norm of an ideal \hyperref[sec:3.5]{III.5}

\medskip
perfect field \hyperref[sec:2.4]{II.4}

polynomial, characteristic \hyperref[sec:2.6]{II.6}

polynomial, minimal \hyperref[sec:2.3]{II.3}

prime field \hyperref[sec:1.7]{I.7}

prime ideal \hyperref[sec:3.3]{III.3}

prime, relatively \hyperref[sec:1.1]{I.1}

primitive element of an extension \hyperref[sec:2.4]{II.4}

primitive root modulo $p$ \hyperref[sec:1.7]{I.7}

primitive root of unity \hyperref[sec:1.6]{I.6}

principal ideal \hyperref[sec:1.1]{I.1}

principal ideal ring \hyperref[sec:1.1]{I.1}

product of ideals \hyperref[sec:3.3]{III.3}

\medskip
quadratic field \hyperref[sec:2.5]{II.5}

quadratic extension \hyperref[sec:6.1]{VI.1}

quasi-algebraically closed field \hyperref[sec:1.7]{I.7}

quaternions \hyperref[sec:5.7]{V.7}

quaternions, Hurwitz \hyperref[sec:5.7]{V.7}

\medskip
ramification, index of \hyperref[sec:5.2]{V.2}

ramified, cf.\ \hyperref[fm:supplement]{the Supplement}

ramify \hyperref[sec:5.3]{V.3}

rank of a module \hyperref[sec:1.5]{I.5}

real quadratic field \hyperref[sec:2.5]{II.5}

reduced ring \hyperref[sec:4.7]{IV.7}

relatively prime elements \hyperref[sec:1.1]{I.1}

remains prime, a prime number or ideal \hyperref[sec:5.4]{V.4}

represent zero \hyperref[sec:1.7]{I.7}

residual degree \hyperref[sec:5.2]{V.2}

residue, quadratic \hyperref[sec:5.4]{V.4}

ring, Dedekind \hyperref[sec:3.4]{III.4}

ring of fractions \hyperref[sec:5.1]{V.1}

ring, Noetherian \hyperref[sec:3.1]{III.1}

ring, reduced \hyperref[sec:4.7]{IV.7}

root of unity, primitive \hyperref[sec:1.6]{I.6}

root, primitive modulo $p$ \hyperref[sec:1.7]{I.7}

\medskip
splitting of a prime number \hyperref[sec:5.4]{V.4}

square-free integer \hyperref[sec:2.5]{II.5}

sum, Gaussian \hyperref[sec:5.5]{V.5}

symbol, Legendre \hyperref[sec:5.5]{V.5}

\medskip
torsion-free module \hyperref[sec:1.5]{I.5}

trace \hyperref[sec:2.6]{II.6}

transcendental element \hyperref[sec:2.3]{II.3}

\medskip
units of a number field \hyperref[sec:4.4]{IV.4}

units of a ring \hyperref[sec:1.1]{I.1}

\medskip
volume of a lattice \hyperref[sec:4.1]{IV.1}

\end{multicols}
